% Part: many-valued-logic
% Chapter: syntax-and-semantics
% Section: formulas

\documentclass[../../../include/open-logic-section]{subfiles}

\begin{document}

\olfileid{mvl}{syn}{fml}

\olsection{\usetoken{P}{formula}}

\begin{defn}[सूत्र]
\ollabel{defn:formulas}
विधानीय भाषा~$\Lang L$ मधील \emph{!!{formula}s} चा संच~$\Frm[L]$
पुढीलप्रमाणे विगमनाने परिभाषित करतो:
\begin{enumerate}
\item प्रत्येक !!{propositional variable}~$\Obj p_i$ हे आण्विक
  !!{formula} असते.
\item $\Lang L$ मधील प्रत्येक $0$-स्थानी संयोजक (विधानीय स्थिर)
  हे आण्विक !!{formula} असते.
\item $\star$ हा $\Lang L$ मधील $n$-स्थानी संयोजक असेल आणि
  $!A_1$, \dots, $!A_n$ ही !!{formula}s असतील, तर
  $\star(!A_1, \dots, !A_n)$ हे !!a{formula} असते.
\tagitem{limitClause}{यांखेरीज काहीही !!a{formula} नाही.}{}
\end{enumerate}
$\star$ $1$-स्थानी असेल, तर $\star(!A_1)$ हे अनेकदा फक्त
$\star !A_1$ असे लिहितात. $\star$ $2$-स्थानी असेल, तर
$\star(!A_1,!A_2)$ हे अनेकदा $(!A_1 \star !A_2)$ असे लिहितात.
\end{defn}

नेहमीप्रमाणे, सर्वांत बाहेरचे कंस आपण अनेकदा न सांगता वगळू.

\begin{ex}
  मानक भाषेत~$\Lang{L_0}$,
  $\Obj p_1 \lif (\Obj p_1 \land \lnot \Obj p_2)$ हे सूत्र आहे.
  गुणाकार तर्कशास्त्राच्या भाषेत ते याऐवजी
  $\Obj p_1 \lif (\Obj p_1 \odot \lnot \Obj p_2)$ असे लिहितात.
  भाषेत $1$-स्थानी $\triangle$ जोडल्यास,
  $\triangle (\Obj p_1 \land \Obj p_2) \lif (\triangle \Obj p_1 \land \triangle \Obj p_2)$
  यांसारखी सूत्रेही मिळतात.
\end{ex}

\end{document}
